\scanpage{66}
\begingroup\makeatletter\tagsleft@true\makeatother
\uline{Fonction de Legendre $\mathrm{sn}\,z$}

La forme $\displaystyle\int\frac{dZ}{\sqrt{(1-Z^2)(1-k^2 Z^2)}}$ est tentante, car elle paraît très symétrique, et rappelle des intégrales élémentaires \unsure{liées|liés} aux fonctions \unsure{trigonométrique|trigonométriques} et \unsure{hyperbolique|hyperboliques} ($k=0,\ k=1$).

Appelons $\mathrm{sn}\,z$ l'intégrale de l'équation
\[
\left(\frac{dZ}{dz}\right)^2=(1-Z^2)(1-k^2Z^2)
\]
qui pd la valeur zéro au point $z=0$, avec $Z'(0)=+1$. L'unicité de cette solution se montre de diverses façons.

1\textsuperscript{o} En faisant un changement de variable convenable, on obtient
\[
Z=\frac{12\unsure{p|\wp}(z-z_0)-(5-k^2)}{12p(z-z_0)-(5k^2-1)}
\qquad
\begin{aligned}
e_1&=\tfrac16(1\unsure{+|\cdot}k^2)\\
e_2&=-\tfrac1{12}(k^2+6k+1)\\
e_3&=-\tfrac1{12}(k^2-6k+1)
\end{aligned}
\]

L'intégrale particulière est
\[
\mathrm{sn}\,z=\frac{12p(z-z_0)-(5-k^2)}{12p(z-z_0)-(5k^2-1)}
\qquad\text{avec}\quad pz_0=\frac{5-k^2}{12}\quad p'z_0=-\frac{1-k^2}{2}\ .\tag*{(1)}
\]

Ces formules sont de forme trop lourde, eu égard à d'autres plus simples. Mais on voit l'existence de $\mathrm{sn}\,z$, qui est une fonction elliptique dont les périodes s'expriment par des intégrales.

\endgroup
% archive sheet number: 65
